\section{Climate Changes: Ongoing and Projected for the Future}
The Arctic is experiencing a phenomenon called Arctic amplification, which is the enhanced warming of the Arctic compared to the rest of the world \citep{serreze2014}. Due to the increasing decline of sea ice during the summer months, the albedo of the Arctic region is decreasing, leading to greater absorption of solar radiation during the time when most solar radiation is available, and thus amplifies the heating. Furthermore, there is a positive local feedback from the lapse rate, as the atmosphere aloft warms slower than the surface due to the polar inversion, which amplifies surface heating \citep{robyn2021}. A study by \citet{rantanen2022} showed that during 1979--2021 the Arctic warmed by more than four times the global average. According to \citet{IPCC2023}, the Arctic is projected to continue warming at a faster rate than the global average. 

Another important aspect of climate change in the Arctic is the thawing of permafrost \citep{serreze2014}. The thawing of permafrost is another positive feedback mechanism that can lead to further warming in the Arctic due to the release of greenhouse gases into the atmosphere, which can further enhance the greenhouse effect \citep{IPCC2023}. The thawing of permafrost can also lead to changes in the landscape, such as the formation of thermokarst (terrain collapse caused by thawing) and thus collapse of infrastructure built on this permafrost. Furthermore, according to \citet{serreze2014} rivers in the Arctic, such as the Mackenzie, Lena, Yenisey, and Ob, may change their flow patterns due to changes in the landscape. This would have significant impacts on the freshwater flow into the Arctic Ocean and sea-ice formation. These changes could significantly alter the heat balance of the Arctic, either preventing the formation of sea ice or, in extreme cases, leading to the Arctic Ocean freezing entirely to the seafloor during winter. Thus, a complete change in the heat balance and important ocean currents from the region may be expected \citep{hartmann2016}. As a consequence, the Arctic is estimated to be ice-free during the summer months by the year 2050 \citep{hartmann2016}. Communities dependent on sea ice for their livelihoods, such as Indigenous peoples, may therefore face even greater challenges as sea ice continues to decline \citep{IPCC2023}.